\begin{table}[htbp]\centering\small\singlespacing
\caption{Directional PVAR external-consistency tests}\label{tab:selected_external_consistency}
\setlength{\tabcolsep}{4pt}
\begin{tabular}{@{}p{4.25cm}crp{6.0cm}@{}}\toprule
Relationship & Expected & Lag-test $p$ & Key response (90\% interval)\\\midrule
Inflation risk $\rightarrow$ headline inflation & $+$ & 0.0299 & M6: $+0.531$ pp $[0.287,\ 0.825]$\\
Depreciation $\rightarrow$ AERI & $+$ & 0.0142 & M2: $+0.0166$ SD $[0.0060,\ 0.0241]$\\
GDP growth $\rightarrow$ AERI & $-$ & 0.0624 & M1: $-0.181$ SD $[-0.338,\ -0.068]$\\
\bottomrule\end{tabular}
\par\medskip\raggedright\footnotesize\textit{Notes:} Expected signs are implied by the economic interpretation of the measures rather than chosen from the estimated responses. All three joint lag tests are statistically significant at the paper's 10 percent criterion, and the reported pointwise 90 percent intervals exclude zero in the expected direction. The inflation-risk and depreciation lag tests are also significant at 5 percent. Exact $p$-values are shown so readers can apply alternative significance conventions. These are conditional dynamic associations under the stated recursive ordering; structural identification is examined separately with the external monetary-policy shock design.
\end{table}
